% Example document using qslides class — comprehensive guide
% Build from package root: make qslides  or  latexmk -pdf tex-qtex/examples/qslides/main.tex
% SPDX-FileCopyrightText: 2026 Frank Langbein <frank@langbein.org>
% SPDX-License-Identifier: LPPL-1.3c
%
% This example serves as a compact guide to qslides-specific features and
% standard beamer capabilities rendered with the Qyber theme.
%
% Class options used here: sans, numeric, titlesplit.
% Other useful options (try uncommenting):
%   Font:        serif, hacker
%   Title page:  titlebasic (default), titlesplit
%   Size:        large (17pt body; bars scale)
%   Citation:    harvard (author-year)
%   Other:       italic, obfuscate (LuaLaTeX blind review)
%   Beamer pass: aspectratio=43, handout, notes
%TC:subst \string\ STR

\documentclass[sans,numeric,titlesplit]{qslides}

\title{Qyber Slides --- Compact Guide}
\subtitle{qslides class \& Beamer features}
\titlebox{%
  \centerline{\includegraphics[width=.5\linewidth]{example-image-a}}
  \smallskip
  The \texttt{qslides} class extends Beamer with the Qyber dark theme,
  custom blocks, layout helpers, and footer-based citations.
}
\author{Author Name}
\date{\today}
\copyrightholder{Frank Langbein}
\licenseyear{2026}
\license{LPPL-1.3c}
\slidesurls{\url{https://qyber.black/}\\\url{https://langbein.org/}}
\setqyberauthorurl{https://langbein.org/}
\slidesbib{refs}

\begin{document}

% ==========================================================================
%  TITLE AND OVERVIEW
% ==========================================================================

\begin{frame}
  \titlepage
\end{frame}

\section{Overview}

\begin{frame}{Overview}
  \begin{columns}[onlytextwidth,T]
    \begin{column}{0.48\linewidth}
      \tableofcontents[sections={1-9}]
    \end{column}
    \begin{column}{0.48\linewidth}
      \tableofcontents[sections={10-17}]
    \end{column}
  \end{columns}
\end{frame}

% ==========================================================================
%  PART I — QSLIDES-SPECIFIC FEATURES
% ==========================================================================

\section{Title Page}

\begin{frame}{Title page variants}
  \framesubtitle{Class options \texttt{titlebasic} (default) and \texttt{titlesplit}}
  \begin{twocolumns}[t]
    \begin{leftcolumn}
      \begin{block}{\texttt{titlebasic}}
        Red header bar with title/subtitle, centred author and date below.
        Optional \texttt{\string\titlegraphic} above the bar.
      \end{block}
    \end{leftcolumn}
    \begin{rightcolumn}
      \begin{block}{\texttt{titlesplit}}
        Red header bar, then a two-column split: left for
        \texttt{\string\titlebox\{...\}} content (image, highlights),
        right for author and date.
      \end{block}
    \end{rightcolumn}
  \end{twocolumns}
  \vspace{1ex}
  \begin{codeblock}{Preamble for split title}
    \texttt{\string\documentclass[titlesplit]\{qslides\}}\\
    \texttt{\string\titlebox\{\string\includegraphics[width=.5\string\linewidth]\{img\}  Some text.\}}
  \end{codeblock}
\end{frame}

\begin{frame}{License footer on title page}
  \framesubtitle{\texttt{\string\license}, \texttt{\string\copyrightholder}, \texttt{\string\licenseyear}}
  When \texttt{\string\license\{SPDX-ID\}} is set, a thin footer bar appears at the
  bottom of the title page with copyright and license information.

  Supported SPDX identifiers include MIT, Apache-2.0, GPL-3.0-only, CC-BY-4.0,
  LPPL-1.3c, and many more (see \texttt{qtex.sty}).
  \begin{codeblock}{Example}
    \texttt{\string\license\{LPPL-1.3c\}}\\
    \texttt{\string\copyrightholder\{Frank Langbein\}}\\
    \texttt{\string\licenseyear\{2026\}}
  \end{codeblock}
\end{frame}

\section{Preamble Commands}

\begin{frame}{Preamble commands}
  \framesubtitle{URLs, author link, bibliography}
  \begin{block}{\texttt{\string\slidesurls\{\string\url\{...\}\}}}
    Sets the URLs displayed in the header/footer right-hand box. Separate
    multiple URLs with \texttt{\string\\}.
  \end{block}
  \begin{block}{\texttt{\string\setqyberauthorurl\{https://...\}}}
    Makes the short author name in the header bar a clickable link.
  \end{block}
  \begin{block}{\texttt{\string\slidesbib\{refs\}}}
    Registers a \texttt{.bib} file.  Citations via
    \texttt{\string\citep}/\texttt{\string\citet} show the full reference
    in a footer bar (via \texttt{\string\bibentry}).
  \end{block}
\end{frame}

\section{Emphasis \& Colours}

\begin{frame}{Emphasis and text colours}
  \framesubtitle{\texttt{\string\emph}, \texttt{\string\hi}, \texttt{\string\alert}}
  \begin{itemize}
    \item \emph{Emphasis} uses \texttt{fg\_emph\_col} (green).
      Syntax: \texttt{\string\emph\{text\}}.
    \item \hi{Highlight} uses \texttt{fg\_highlight\_col} as background
      with \texttt{fg\_body\_col} (white) text.
      Syntax: \texttt{\string\hi\{text\}}.
    \item \alert{Alert} uses \texttt{fg\_highlight\_col} for the text colour.
      Syntax: \texttt{\string\alert\{text\}} (standard Beamer).
  \end{itemize}
  \vspace{1ex}
  Links: \href{https://qyber.black/}{Qyber} uses \texttt{fg\_link\_col} (blue).
  \url{https://langbein.org/} uses \texttt{fg\_ref\_col} (green).
\end{frame}

\begin{frame}{Colour palette --- backgrounds}
  \framesubtitle{Named \texttt{bg\_*\_col} colours used by blocks, bars, and page}
  \begin{columns}[onlytextwidth,T]
    \begin{column}{0.48\linewidth}
      \textbf{Core backgrounds}\\[0.5ex]
      \colorbox{bg_page_col}{\color{fg_body_col}\strut\hspace{0.9em}}\ \texttt{bg\_page\_col}\\[0.3ex]
      \colorbox{bg_surface_col}{\color{fg_body_col}\strut\hspace{0.9em}}\ \texttt{bg\_surface\_col}\\[0.3ex]
      \colorbox{bg_header_col}{\color{fg_body_col}\strut\hspace{0.9em}}\ \texttt{bg\_header\_col}\\[0.3ex]
      \colorbox{bg_highlight_col}{\color{fg_body_col}\strut\hspace{0.9em}}\ \texttt{bg\_highlight\_col}\\[0.3ex]
      \colorbox{bg_video_col}{\color{fg_body_col}\strut\hspace{0.9em}}\ \texttt{bg\_video\_col}
    \end{column}
    \begin{column}{0.48\linewidth}
      \textbf{Block backgrounds}\\[0.5ex]
      \colorbox{bg_info_col}{\color{fg_body_col}\strut\ info\hspace{0.5em}}\ \texttt{bg\_info\_col}\\[0.3ex]
      \colorbox{bg_error_col}{\color{fg_body_col}\strut\ error\hspace{0.5em}}\ \texttt{bg\_error\_col}\\[0.3ex]
      \colorbox{bg_example_col}{\color{fg_body_col}\strut\ example\hspace{0.5em}}\ \texttt{bg\_example\_col}\\[0.3ex]
      \colorbox{bg_code_col}{\color{fg_body_col}\strut\ code\hspace{0.5em}}\ \texttt{bg\_code\_col}\\[0.3ex]
      \colorbox{bg_result_col}{\color{fg_body_col}\strut\ result\hspace{0.5em}}\ \texttt{bg\_result\_col}\\[0.3ex]
      \colorbox{bg_success_col}{\color{fg_body_col}\strut\ success\hspace{0.5em}}\ \texttt{bg\_success\_col}
    \end{column}
  \end{columns}
\end{frame}

\begin{frame}{Colour palette --- foregrounds \& accents}
  \framesubtitle{Named \texttt{fg\_*\_col} and \texttt{ac\_*\_col} colours}
  \begin{columns}[onlytextwidth,T]
    \begin{column}{0.48\linewidth}
      \textbf{Text foregrounds}\\[0.5ex]
      {\color{fg_body_col}\rule{0.9em}{1.2ex}}\ \texttt{fg\_body\_col} (white)\\[0.3ex]
      {\color{fg_head_col}\rule{0.9em}{1.2ex}}\ \texttt{fg\_head\_col} (gold)\\[0.3ex]
      {\color{fg_footer_col}\rule{0.9em}{1.2ex}}\ \texttt{fg\_footer\_col} (grey)\\[0.3ex]
      {\color{fg_highlight_col}\rule{0.9em}{1.2ex}}\ \texttt{fg\_highlight\_col} (orange)\\[0.3ex]
      {\color{fg_emph_col}\rule{0.9em}{1.2ex}}\ \texttt{fg\_emph\_col} (green)\\[0.3ex]
      {\color{fg_header_bar_col}\rule{0.9em}{1.2ex}}\ \texttt{fg\_header\_bar\_col}
    \end{column}
    \begin{column}{0.48\linewidth}
      \textbf{Links}\\[0.5ex]
      {\color{fg_link_col}\rule{0.9em}{1.2ex}}\ \texttt{fg\_link\_col} (blue)\\[0.3ex]
      {\color{fg_ref_col}\rule{0.9em}{1.2ex}}\ \texttt{fg\_ref\_col} (green)\\[1.5ex]
      \textbf{Accents}\\[0.5ex]
      {\color{ac_rule_col}\rule{0.9em}{1.2ex}}\ \texttt{ac\_rule\_col} (bright green)\\[0.3ex]
      {\color{ac_en_symbol_col}\rule{0.9em}{1.2ex}}\ \texttt{ac\_en\_symbol\_col} (purple)
    \end{column}
  \end{columns}
\end{frame}

\section{Blocks}

\begin{frame}{Blocks (6 types, with header)}
  \begin{columns}[onlytextwidth,T]
    \begin{column}{0.48\linewidth}
      \begin{block}{Info block}
        \texttt{block}: definitions (cyan).
      \end{block}
      \begin{alertblock}{Error block}
        \texttt{alertblock}: errors (red).
      \end{alertblock}
      \begin{exampleblock}{Example block}
        \texttt{exampleblock}: examples (violet).
      \end{exampleblock}
    \end{column}
    \begin{column}{0.48\linewidth}
      \begin{codeblock}{Code block}
        \texttt{codeblock}: code (blue).
      \end{codeblock}
      \begin{resultblock}{Result block}
        \texttt{resultblock}: results (olive).
      \end{resultblock}
      \begin{successblock}{Success block}
        \texttt{successblock}: success (green).
      \end{successblock}
    \end{column}
  \end{columns}
\end{frame}

\begin{frame}{Blocks without header}
  \framesubtitle{Empty title: \texttt{\string\begin\{block\}\{\string\}\string\end\{block\}}}
  \begin{columns}[onlytextwidth,T]
    \begin{column}{0.48\linewidth}
      \begin{block}{}
        \texttt{block}\{\}: info (cyan).
      \end{block}
      \begin{alertblock}{}
        \texttt{alertblock}\{\}: error (red).
      \end{alertblock}
      \begin{exampleblock}{}
        \texttt{exampleblock}\{\}: example (violet).
      \end{exampleblock}
    \end{column}
    \begin{column}{0.48\linewidth}
      \begin{codeblock}{}
        \texttt{codeblock}\{\}: code (blue).
      \end{codeblock}
      \begin{resultblock}{}
        \texttt{resultblock}\{\}: result (olive).
      \end{resultblock}
      \begin{successblock}{}
        \texttt{successblock}\{\}: success (green).
      \end{successblock}
    \end{column}
  \end{columns}
\end{frame}

\section{Layout Helpers}

\begin{frame}{Lists}
  \framesubtitle{Itemize and enumerate (3 nesting levels)}
  \begin{columns}[onlytextwidth,T]
    \begin{column}{0.48\linewidth}
      \textbf{Itemize}
      \begin{itemize}
        \item First bullet
          \begin{itemize}
            \item Subitem 1
            \item Subitem 2
          \end{itemize}
        \item Second bullet
        \item Third bullet
          \begin{itemize}
            \item Subitem 1
              \begin{itemize}
                \item Subsubitem 1
                \item Subsubitem 2
              \end{itemize}
            \item Subitem 2
          \end{itemize}
      \end{itemize}
    \end{column}
    \begin{column}{0.48\linewidth}
      \textbf{Enumerate}
      \begin{enumerate}
        \item First step
          \begin{enumerate}
            \item Substep 1
            \item Substep 2
          \end{enumerate}
        \item Second step
        \item Third step
          \begin{enumerate}
            \item Substep 1
              \begin{enumerate}
                \item Subsubstep 1
                \item Subsubstep 2
              \end{enumerate}
          \end{enumerate}
      \end{enumerate}
    \end{column}
  \end{columns}
\end{frame}

\begin{frame}{Headerless boxes (\texttt{slidebox})}
  \framesubtitle{\texttt{slidebox} and \texttt{slidebox[border]}}
  Content-only boxes with no header bar. Optional thin border.
  \begin{slidebox}
    \textbf{No border} --- normal text colours, no semantic colouring.
    Use for grouping related content.
    Syntax: \texttt{slidebox} environment.
  \end{slidebox}
  \begin{slidebox}[border]
    \textbf{With border} --- same style plus a thin rule.
    Useful for call-outs or framing important text.
    Syntax: \texttt{slidebox} with \texttt{[border]} option.
  \end{slidebox}
\end{frame}

\begin{frame}{Two-column slide}
  \framesubtitle{\texttt{twocolumns}, \texttt{leftcolumn}, \texttt{rightcolumn}}
  \begin{twocolumns}[t]
    \begin{leftcolumn}
      \textbf{Left column}
      \begin{itemize}
        \item First point
        \item Second point
        \item Third point
      \end{itemize}
    \end{leftcolumn}
    \begin{rightcolumn}
      \textbf{Right column}
      \begin{itemize}
        \item Another point
        \item Final point
        \item One more
      \end{itemize}
    \end{rightcolumn}
  \end{twocolumns}
  \vspace{1ex}
  Alignment options: \texttt{[t]} top, \texttt{[c]} centre, \texttt{[b]} bottom.
  Each column is 48\% of \texttt{\string\linewidth}.
\end{frame}

\begin{frame}{Two columns as part of a slide}
  \framesubtitle{Blocks inside columns; intro text above}
  Intro text above the columns shows that \texttt{twocolumns} can share
  the slide with other content.
  \begin{twocolumns}[t]
    \begin{leftcolumn}
      \begin{block}{Left block}
        Content in the left column inside a standard info block.
      \end{block}
    \end{leftcolumn}
    \begin{rightcolumn}
      \begin{codeblock}{Right block}
        Content in the right column inside a code block.
      \end{codeblock}
    \end{rightcolumn}
  \end{twocolumns}
\end{frame}

\begin{frame}{Centred column alignment}
  \begin{twocolumns}[c]
    \begin{leftcolumn}
      Short left text.\\
      Another line.\\
      Another line.\\
      Another line.\\
      Another line.
    \end{leftcolumn}
    \begin{rightcolumn}
      Short right text; columns are vertically centred.
    \end{rightcolumn}
  \end{twocolumns}
\end{frame}

\begin{frame}{Placed boxes (relative coordinates)}
  \framesubtitle{\texttt{\string\placebox\{x\}\{y\}\{w\}\{h\}\{content\}}}
  Position and size in $[0,1]$ relative to the text area. $(0,0)$ = bottom-left.
  \placebox{0.05}{0.0}{0.4}{0.3}{%
    \begin{block}{Bottom-left}
      At $(0.05, 0.0)$, size $0.4 \times 0.3$.
  \end{block}}%
  \placebox{0.55}{0.15}{0.4}{0.35}{%
    \begin{codeblock}{Right side}
      At $(0.55, 0.15)$, size $0.4 \times 0.35$.
  \end{codeblock}}%
  \placebox{0.25}{0.55}{0.5}{0.2}{%
    \begin{exampleblock}{Centre}
      Overlapping boxes are drawn in source order.
  \end{exampleblock}}%
\end{frame}

\begin{frame}{Placed box with optional anchor}
  \framesubtitle{\texttt{\string\placebox[anchor]\{x\}\{y\}\{w\}\{h\}\{content\}}}
  The default anchor is \texttt{south west}.  With \texttt{[north east]},
  the given $(x,y)$ is the top-right corner of the box.
  \placebox[north east]{0.95}{0.85}{0.45}{0.25}{%
    \begin{block}{Top-right anchored}
      Anchored \texttt{north east} at $(0.95, 0.85)$.
  \end{block}}%
  \placebox{0.05}{0.05}{0.45}{0.25}{%
    \begin{resultblock}{Bottom-left (default)}
      Anchored \texttt{south west} at $(0.05, 0.05)$.
  \end{resultblock}}%
\end{frame}

\section{Video Layouts}

\begin{frame}{Video layout modes}
  \framesubtitle{Five background templates for video-integrated presentations}
  The Qyber theme provides five layout modes for slides that include a
  video area (displayed as \texttt{bg\_video\_col}).
  \begin{description}
    \item[\texttt{full}] Default; content fills the entire slide.
    \item[\texttt{left}] Content left (60\%), video right (40\%).
    \item[\texttt{right}] Video left (40\%), content right (60\%).
    \item[\texttt{corner}] Small video in one corner (25\%$\times$25\%).
    \item[\texttt{overlay}] Full-slide video background, content overlay box.
  \end{description}
\end{frame}

\begin{frame}{Video layout commands}
  \framesubtitle{Selecting layout, corner side, and video colour}
  \begin{codeblock}{Setting the layout}
    \texttt{\string\setbeamertemplate\{background\}[layoutleft]}\\
    \texttt{\string\setbeamertemplate\{background\}[layoutcorner]}\\
    \texttt{\string\setbeamertemplate\{background\}[layoutfull]}
  \end{codeblock}
  \begin{codeblock}{Corner side and video colour}
    \texttt{\string\setcornerleft}\ \ or\ \ \texttt{\string\setcornerright}\\[0.5ex]
    \texttt{\string\setvideocolor\{0.0,0.2,0.8\}}
  \end{codeblock}
  Reset to normal with \texttt{[layoutfull]} after the video section.
\end{frame}

\begin{frame}{Left split layout}
  \setbeamertemplate{background}[layoutleft]
  \framesubtitle{Content area left (60\%), video area right (40\%)}
  This slide uses \texttt{[layoutleft]}.
  The dark area on the right would contain
  a video feed in a live presentation.
  \begin{itemize}
    \item Content stays in the left 60\%.
    \item The header bar spans the content area only.
    \item Reset with \texttt{[layoutfull]}.
  \end{itemize}
\end{frame}

\setbeamertemplate{background}[layoutfull]

\begin{frame}{Footer bar mode}
  \framesubtitle{\texttt{\string\enablefooter} moves the header bar to the bottom}
  Normally the red title bar is at the top (\texttt{header} mode, default).
  Calling \texttt{\string\enablefooter} switches to \texttt{footer} mode:
  the bar moves to the bottom, freeing the top for full-bleed content.
  \begin{codeblock}{Usage}
    \texttt{\string\enablefooter}\\
    \texttt{\string\begin\{frame\}\{Title\}}\\
    \texttt{\ \ Content above the bottom bar.}\\
    \texttt{\string\end\{frame\}}
  \end{codeblock}
\end{frame}

\section{Class Options}

\begin{frame}{Class options summary}
  \framesubtitle{All options for \texttt{\string\documentclass[...]\{qslides\}}}
  \small
  \begin{columns}[onlytextwidth,T]
    \begin{column}{0.48\linewidth}
      \begin{block}{Font}
        \texttt{sans} (default), \texttt{serif}, \texttt{hacker}
      \end{block}
      \begin{block}{Title page}
        \texttt{titlebasic} (default), \texttt{titlesplit}
      \end{block}
      \begin{block}{Size}
        12pt default; \texttt{large} = 17pt
      \end{block}
    \end{column}
    \begin{column}{0.48\linewidth}
      \begin{block}{Citation}
        \texttt{numeric} (default), \texttt{harvard}
      \end{block}
      \begin{block}{Other qtex options}
        \texttt{italic}, \texttt{obfuscate},\\
        \texttt{parskip}/\texttt{parindent},\\
        \texttt{single}/\texttt{onehalfspacing}
      \end{block}
      \begin{block}{Beamer pass-through}
        \texttt{aspectratio=43}, \texttt{handout},\\
        \texttt{notes}, any Beamer option.
      \end{block}
    \end{column}
  \end{columns}
\end{frame}

\section{Footnotes \& Citations}

\begin{frame}{Footnotes}
  \framesubtitle{\texttt{\string\footnote\{...\}} renders in a TikZ-overlay footer bar}
  Footnotes appear as \texttt{[N]} markers in the text\footnote{This is the
  first footnote. It appears at the bottom in a footer bar.} and the full
  text is collected into a bar at the slide bottom, matching the title-page
  license bar style.

  Multiple footnotes\footnote{Second footnote, demonstrating numbering.}
  are numbered per slide; the counter resets on each new frame.
\end{frame}

\begin{frame}{Citations and bibliography}
  \framesubtitle{\texttt{\string\citep} and \texttt{\string\citet} with footer references}
  With \texttt{\string\slidesbib\{refs\}} in the preamble,
  \texttt{\string\citep} and \texttt{\string\citet} show the numeric marker
  in the text and put the full reference in the footer bar via
  \texttt{\string\bibentry}.

  Example: \citep{langbein2026} describes the theme,
  while \citet{knuth1984} is the classic \TeX{} reference.
  \begin{block}{Setup}
    \small
    \texttt{\string\slidesbib\{refs\}} in preamble registers the \texttt{.bib} file.\\
    References appear only in the footer; there is no bibliography slide.
  \end{block}
\end{frame}

% ==========================================================================
%  PART II — STANDARD BEAMER FEATURES (Qyber theme)
% ==========================================================================

\section{Overlays}

\begin{frame}{Overlays: \texttt{\string\pause}}
  \framesubtitle{Content appears step by step}
  \begin{itemize}
    \item This is visible from the start.
      \pause
    \item This appears on the second step.
      \pause
    \item This on the third.
  \end{itemize}
  \vspace{1ex}
  After the final pause, all items are visible.
\end{frame}

\begin{frame}{Overlays: \texttt{\string\only} and \texttt{\string\onslide}}
  \framesubtitle{Fine-grained overlay control}
  \begin{itemize}
    \item \texttt{\string\only<2>\{text\}}: \only<2>{Visible on slide 2 only.}
    \item \texttt{\string\onslide<3->\{text\}}: \onslide<3->{Visible from slide 3 onwards.}
    \item \texttt{\string\uncover<2-3>\{text\}}: \uncover<2-3>{Visible on slides 2 and 3 (space reserved).}
    \item \texttt{\string\visible<4>\{text\}}: \visible<4>{Visible on slide 4 only (space reserved).}
  \end{itemize}
  \vspace{1ex}
  \only<1>{Slide 1: only the first item and this text.}%
  \only<2>{Slide 2: \texttt{\string\only} and \texttt{\string\uncover} items appear.}%
  \only<3>{Slide 3: \texttt{\string\onslide} and \texttt{\string\uncover} items shown.}%
  \only<4>{Slide 4: \texttt{\string\visible} item shown.}%
\end{frame}

\begin{frame}{Item-level overlay specifications}
  \framesubtitle{\texttt{<n->} on list items}
  \begin{enumerate}
    \item<1-> Visible from slide 1
    \item<2-> Visible from slide 2
    \item<3-> Visible from slide 3
    \item<1,3> Visible on slides 1 and 3 only
  \end{enumerate}
  \vspace{1ex}
  \texttt{\string\alt} example:
  \alt<2>{This text is shown on slide 2.}{This text is shown on all other slides.}
\end{frame}

\begin{frame}{Overlay actions on text}
  \framesubtitle{\texttt{\string\textbf<n>}, \texttt{\string\color<n>}, \texttt{\string\action<spec>}}
  \begin{itemize}
    \item Normal text, \textbf<2>{bold on slide 2}, normal again.
    \item \action<3| alert@3>{Alerted on slide 3.}
    \item \temporal<2>{Before slide 2}{On slide 2}{After slide 2}
  \end{itemize}
\end{frame}

\section{Columns}

\begin{frame}{Beamer columns (raw)}
  \framesubtitle{\texttt{columns} / \texttt{column} environments}
  \begin{columns}[T]
    \begin{column}{0.32\textwidth}
      \begin{block}{Column 1}
        One-third width. Beamer's native \texttt{columns} environment
        gives full control over widths.
      \end{block}
    \end{column}
    \begin{column}{0.32\textwidth}
      \begin{exampleblock}{Column 2}
        Another third.  Any block type works inside columns.
      \end{exampleblock}
    \end{column}
    \begin{column}{0.32\textwidth}
      \begin{alertblock}{Column 3}
        Final third. Use \texttt{[T]} for top alignment or
        \texttt{[c]} for centred.
      \end{alertblock}
    \end{column}
  \end{columns}
\end{frame}

\section{Mathematics}

\begin{frame}{Mathematics}
  \framesubtitle{Inline, display, and numbered equations}
  Inline: $a^2 + b^2 = c^2$, and $\sum_{i=1}^n i = \frac{n(n+1)}{2}$.
  \begin{displaymath}
    \int_0^{+\infty} e^{-x^2}\,\mathrm{d}x = \frac{\sqrt{\pi}}{2}
  \end{displaymath}
  Numbered:
  \begin{equation}
    \nabla \cdot \mathbf{E} = \frac{\rho}{\varepsilon_0}
  \end{equation}
  Aligned:
  \begin{align}
    f(x) &= x^2 + 2x + 1 \\
    &= (x+1)^2
  \end{align}
\end{frame}

\begin{frame}{Theorem environments}
  \framesubtitle{Beamer's built-in theorem, proof, definition, example}
  \begin{theorem}[Pythagoras]
    In a right triangle with legs $a$, $b$ and hypotenuse $c$:
    $a^2 + b^2 = c^2$.
  \end{theorem}
  \begin{proof}
    By area comparison of squares on each side. \qedhere
  \end{proof}
  \begin{definition}
    A \emph{prime number} is a natural number greater than 1 that has no
    positive divisors other than 1 and itself.
  \end{definition}
\end{frame}

\section{Figures \& Tables}

\begin{frame}{Figures and tables}
  \framesubtitle{Graphics and tabular side by side}
  \begin{columns}[onlytextwidth,c]
    \begin{column}{0.48\linewidth}
      \centering
      \includegraphics[width=0.8\linewidth]{example-image-b}
      \captionof{figure}{A placeholder image.}
    \end{column}
    \begin{column}{0.48\linewidth}
      \centering
      \begin{tabular}{lcc}
        \hline
        Method & Accuracy & Time \\
        \hline
        Baseline & 78.2\% & 1.2s \\
        Ours     & 91.5\% & 0.8s \\
        Oracle   & 99.1\% & 3.4s \\
        \hline
      \end{tabular}
      \captionof{table}{Comparison of methods.}
    \end{column}
  \end{columns}
\end{frame}

\section{Hyperlinks \& Navigation}

\begin{frame}{Hyperlinks and navigation buttons}
  \framesubtitle{\texttt{\string\hyperlink}, \texttt{\string\beamerbutton}, etc.}
  \hypertarget{linktarget}{}%
  \begin{itemize}
    \item \texttt{\string\hyperlink\{label\}\{text\}}: jump to a labelled frame.
    \item \texttt{\string\beamerbutton\{text\}}: \beamerbutton{Button}
    \item \texttt{\string\beamergotobutton\{text\}}: \beamergotobutton{Go to}
    \item \texttt{\string\beamerreturnbutton\{text\}}: \beamerreturnbutton{Return}
    \item \texttt{\string\beamerskipbutton\{text\}}: \beamerskipbutton{Skip}
  \end{itemize}
  \vspace{1ex}
  Example link: \hyperlink{linktarget}{\beamergotobutton{Jump to this slide}}.
\end{frame}

\section{Frame Options}

\begin{frame}[plain]{}
  \vfill
  \begin{center}
    {\Large\textcolor{fg_head_col}{Plain frame}}\\[2ex]
    \texttt{[plain]} removes headers and footers.\\
    Useful for full-bleed images or section dividers.
  \end{center}
  \vfill
\end{frame}

\begin{frame}[fragile]{Fragile frame}
  \framesubtitle{\texttt{[fragile]} required for verbatim content}
  Use \texttt{[fragile]} when the frame contains verbatim or code:
  \begin{semiverbatim}
    \\documentclass[sans,numeric]\{qslides\}
    \\title\{My Talk\}
    \\author\{A. Speaker\}
    \\begin\{document\}
    \\begin\{frame\}\\titlepage\\end\{frame\}
    \\begin\{frame\}\{Hello\}Content.\\end\{frame\}
    \\end\{document\}
  \end{semiverbatim}
\end{frame}

\begin{frame}{Frame alignment options}
  \framesubtitle{\texttt{[t]} (default), \texttt{[c]}, \texttt{[b]}}
  \begin{twocolumns}[t]
    \begin{leftcolumn}
      \begin{block}{\texttt{[t]} top (default in qslides)}
        Content starts at the top of the frame body.
        This is the qslides default because \texttt{LoadClass[...,t]\{beamer\}}
        is set in the class file.
      \end{block}
    \end{leftcolumn}
    \begin{rightcolumn}
      \begin{block}{\texttt{[c]} and \texttt{[b]}}
        Add \texttt{[c]} to a frame for vertically centred content, or
        \texttt{[b]} for bottom-aligned content.  These override the class
        default for that frame only.
      \end{block}
    \end{rightcolumn}
  \end{twocolumns}
\end{frame}

\begin{frame}[label=labelexample]{Labelled frame}
  \framesubtitle{\texttt{[label=name]} for cross-references}
  This frame has \texttt{[label=labelexample]}.  Reference it with\\
  \texttt{\string\hyperlink\{labelexample\}\{...\}}
  or use it as an \texttt{\string\againframe} target.
  \begin{block}{Other frame options}
    \small
    \begin{description}
      \item[\texttt{allowframebreaks}] Auto-split long content across slides (useful for bibliographies).
      \item[\texttt{shrink}] Shrink content to fit; use sparingly.
      \item[\texttt{squeeze}] Reduce vertical spacing to fit more.
    \end{description}
  \end{block}
\end{frame}

\section{Sections \& TOC}

\begin{frame}{Sections and navigation}
  \framesubtitle{\texttt{\string\section}, \texttt{\string\subsection}, \texttt{\string\tableofcontents}}
  \begin{twocolumns}[t]
    \begin{leftcolumn}
      \begin{itemize}
        \item \texttt{\string\section\{Name\}}
        \item \texttt{\string\subsection\{Name\}}
        \item \texttt{\string\tableofcontents}
        \item \texttt{[currentsection]} highlights active section
      \end{itemize}
    \end{leftcolumn}
    \begin{rightcolumn}
      \begin{codeblock}{Auto-outline}
        \small
        \texttt{\string\AtBeginSection[]\{}\\
          \texttt{\ \ \string\begin\{frame\}\{Outline\}}\\
          \texttt{\ \ \ \ \string\tableofcontents[currentsection]}\\
        \texttt{\ \ \string\end\{frame\}\}}
      \end{codeblock}
    \end{rightcolumn}
  \end{twocolumns}
\end{frame}

\begin{frame}{Current section highlight}
  \begin{columns}[onlytextwidth,T]
    \begin{column}{0.48\linewidth}
      \tableofcontents[sections={1-9},currentsection]
    \end{column}
    \begin{column}{0.48\linewidth}
      \tableofcontents[sections={10-17},currentsection]
    \end{column}
  \end{columns}
\end{frame}

\section{Verbatim \& Code}

\begin{frame}[fragile]{Verbatim and inline code}
  \framesubtitle{\texttt{verbatim} environment and \texttt{\string\verb}}
  Inline: \verb|\documentclass{qslides}| renders in monospace.

  Block verbatim:
  \begin{semiverbatim}
    \\begin\{frame\}\{Title\}
    \\begin\{block\}\{Info\}
    Content here.
    \\end\{block\}
    \\end\{frame\}
  \end{semiverbatim}
\end{frame}

\section{Multimedia}

\begin{frame}{Multimedia (overview)}
  \framesubtitle{\texttt{\string\movie}, \texttt{\string\sound}, \texttt{\string\animate}}
  Beamer supports embedding multimedia via the \texttt{multimedia} package:
  \begin{block}{Commands}
    \begin{itemize}
      \item \texttt{\string\movie\{poster\}\{file\}} --- embed a video.
        Options: \texttt{autostart}, \texttt{loop}, \texttt{width}, \texttt{height}.
      \item \texttt{\string\sound\{text\}\{file\}} --- embed audio.
      \item \texttt{\string\animate} --- animate overlays.
    \end{itemize}
  \end{block}
  \begin{alertblock}{Note}
    Multimedia support depends on the PDF viewer. Tested with Okular and
    Adobe Acrobat. External files must be present at build time.
  \end{alertblock}
\end{frame}

\section{Notes}

\begin{frame}{Speaker notes}
  \framesubtitle{\texttt{\string\note\{...\}} and presentation modes}
  \note{This is a speaker note visible in notes mode.}
  \begin{itemize}
    \item \texttt{\string\note\{text\}} attaches a note to the current slide.
    \item \texttt{\string\note[item]\{text\}} adds to a note list.
    \item Build with \texttt{notes} option to show notes:
      \texttt{\string\documentclass[notes]\{qslides\}}.
    \item \texttt{handout} mode suppresses overlays for print.
  \end{itemize}
  \begin{codeblock}{Dual-screen setup}
    \small
    \texttt{\string\documentclass[notes]\{qslides\}}\\
    \texttt{\string\setbeameroption\{show notes on second screen=right\}}
  \end{codeblock}
\end{frame}

\section{Appendix}

\begin{frame}{Appendix slides}
  \framesubtitle{\texttt{\string\appendix} and backup slides}
  \begin{block}{Usage}
    Place \texttt{\string\appendix} after the last regular slide.
    Slides after \texttt{\string\appendix} do not count towards the
    total frame number shown in the header.
  \end{block}
  \begin{codeblock}{Convention}
    \small
    \texttt{\string\appendix}\\
    \texttt{\string\begin\{frame\}\{Backup: Extra Data\}}\\
    \texttt{\ \ Additional material for Q\&A.}\\
    \texttt{\string\end\{frame\}}
  \end{codeblock}
\end{frame}

\section{Transitions}

\begin{frame}{Slide transitions (overview)}
  \framesubtitle{PDF viewer dependent}
  Beamer can embed slide transition effects into the PDF.  These require
  a supporting viewer (e.g.\ Okular, Adobe Acrobat in full-screen mode).
  \begin{block}{Commands}
    \small
    \begin{columns}[onlytextwidth,T]
      \begin{column}{0.48\linewidth}
        \begin{itemize}
          \item \texttt{\string\transdissolve}
          \item \texttt{\string\transblindshorizontal}
          \item \texttt{\string\transblindsvertical}
          \item \texttt{\string\transboxin}
        \end{itemize}
      \end{column}
      \begin{column}{0.48\linewidth}
        \begin{itemize}
          \item \texttt{\string\transboxout}
          \item \texttt{\string\transwipe}
          \item \texttt{\string\transcover}
          \item \texttt{\string\transfade}
        \end{itemize}
      \end{column}
    \end{columns}
  \end{block}
  Apply to a single frame or globally with\\
  {\small\texttt{\string\addtobeamertemplate\{background\}\{\string\transfade\}\{\}}}.
\end{frame}

\section{Description Lists}

\begin{frame}{Description lists}
  \framesubtitle{Beamer \texttt{description} environment}
  \begin{description}
    \item[qslides] Beamer-based slide class with Qyber dark theme.
    \item[qarticle] Article-style class for papers (no chapters).
    \item[qbook] Book-style class for theses and long documents.
    \item[qtex.sty] Shared package: fonts, layout, citations, metadata.
  \end{description}
\end{frame}

\section{Summary}

\begin{frame}{Summary}
  \framesubtitle{Quick reference card}
  \small\vskip-2ex%
  \begin{columns}[onlytextwidth,T]
    \begin{column}{0.48\linewidth}
      \begin{block}{Qslides features}
        \begin{description}
          \item[Title] \texttt{titlebasic}/\texttt{titlesplit}, \texttt{\string\titlebox}
          \item[Blocks] 6 types $\times$ 2 (with/without header)
          \item[Layout] \texttt{slidebox}, \texttt{twocolumns}, \texttt{placebox}
          \item[Video] 5 layout modes + \texttt{\string\enablefooter}
          \item[Cite] \texttt{\string\citep}/\texttt{\string\citet} in footer
          \item[Colours] Full named palette (\texttt{bg\_*}, \texttt{fg\_*}, \texttt{ac\_*})
        \end{description}
      \end{block}
    \end{column}
    \begin{column}{0.48\linewidth}
      \begin{block}{Beamer features}
        \begin{description}
          \item[Overlays] \texttt{\string\pause}, \texttt{\textless n-\textgreater}, \texttt{\string\only}, \texttt{\string\alt}
          \item[Maths] Equations, theorems, proofs
          \item[Media] \texttt{\string\movie}, \texttt{\string\sound}
          \item[Frames] \texttt{plain}, \texttt{fragile}, \texttt{label}
          \item[Nav] Sections, TOC, hyperlink buttons
          \item[Notes] \texttt{\string\note}, handout mode
        \end{description}
      \end{block}
    \end{column}
  \end{columns}
\end{frame}

\end{document}
